\scanpage{37}
On peut songer à la version additive, en la définissant d'abord sur les $E_{\unsure{\eta|y}}$ simples (i.e. abs. simples, quitte à permettre des extensions finies sur $\mathbf{Q}_{\unsure{\ell|l}}$) et en prolongeant par linéarité, ce qui revient à définir les facteurs locaux d'un $L_{E_{\unsure{\eta|y}}}(\unsure{t})$ par $L_{E_{\unsure{\eta|y}}(x)}(\unsure{t})$, en posant $E_{\unsure{\eta|y}}(x)=\sum E_{\unsure{\eta|y}}^{i*}(x)$, les $E_{\unsure{\eta|y}}^i$ étant les facteurs simples dans une suite de composition de $E$. Ceci n'est pas encore satisfaisant, car les $E_{\unsure{\eta|y}}(x)$ ainsi construits ne sont pas de \unsure{construction} locale en $(E,x)$, -- \unsure{p. ex.} ne commutent pas à l'extension \unsure{étale|ét.} \deleted{sur $X$} (à ext. résiduelle \inserted{en $x$} \deleted{en} triviale) sur $X$.

\margin{36}

